% Generated draft: requires mathematical and source review.
\documentclass[11pt]{article}
\usepackage[margin=25mm]{geometry}
\usepackage{fontspec}
\setmainfont{Latin Modern Roman}
\usepackage{amsmath,amssymb,mathtools,unicode-math}
\setmathfont{Latin Modern Math}
\usepackage{enumitem,needspace,xurl,hyperref}
\setlength{\parindent}{0pt}
\setlength{\parskip}{5pt}
\setlength{\emergencystretch}{3em}
\newcommand{\N}{\mathbb{N}}
\newcommand{\Z}{\mathbb{Z}}
\newcommand{\Q}{\mathbb{Q}}
\newcommand{\R}{\mathbb{R}}
\newcommand{\C}{\mathbb{C}}
\newcommand{\F}{\mathbb{F}}
\newcommand{\A}{\mathbb{A}}
\newcommand{\PP}{\mathbb{P}}
\newcommand{\E}{\mathbb{E}}
\newcommand{\eps}{\varepsilon}
\newcommand{\dd}{\,\mathrm{d}}
\begin{document}
\section*{UnsolvedMath ID: 30005406}
\textbf{Source problem code:} catalogue entry
\subsection{Definitions and mathematical statement}
Let \(p\) be a prime and let \(k\) be a field of characteristic \(p>0\). Let \(X\) be a flag variety over \(k\), and let \(\operatorname{Pic}(X)\) denote the group of isomorphism classes of line bundles on \(X\). A line bundle is a locally free \(\mathcal O_X\)-module of rank one. For a line bundle \(\mathcal L\) and an integer \(i\geq 0\), write \(H^i(X,\mathcal L)\) for its sheaf cohomology, regarded as a \(k\)-vector space.

\textbf{Problem (statement scheme with unresolved scope).} For every positive characteristic and every field and flag variety within the intended scope of the source, determine
\[
H^i(X,\mathcal L)
\qquad
\text{for every }[\mathcal L]\in\operatorname{Pic}(X)
\text{ and every integer }i\geq 0.
\]
In particular, give necessary and sufficient conditions on the parameters \(k,X,\mathcal L,i\) for
\[
H^i(X,\mathcal L)=0,
\qquad\text{equivalently, characterize when }H^i(X,\mathcal L)\neq 0.
\]

\textbf{Unresolved qualifications.} The supplied record does not specify whether \(k\) is algebraically closed, whether the flag varieties are complete or partial, or which algebraic groups and splitting hypotheses are intended. In the usual homogeneous-space description \(X=G/P\), with \(G\) a connected reductive algebraic group and \(P\) a parabolic subgroup, these choices require clarification; this description does not settle the intended scope. The record also does not specify whether determining cohomology means determining only the underlying vector spaces or their natural group representations when such structures are available. Consequently, the displayed request is a statement scheme, not a resolution of these missing hypotheses.

The supplied original excerpt attributes a characterization in degree \(0\) to Kempf's vanishing theorem, with source-local references [9, 8], and a characterization in degree \(1\) to Andersen, with source-local reference [2]. No precise versions or hypotheses of those results are supplied. The excerpt ends mid-sentence after ``but general results'', followed by a report page artifact; its missing continuation is not reconstructed here.
\subsection{Short English statement}
Determine the cohomology of every line bundle on the flag varieties intended by the source in characteristic \(p>0\), including exact criteria for vanishing in each degree. The allowed fields, flag varieties, and desired representation structure remain unspecified. The supplied dated curation labels the broad problem partially solved and reports substantial special cases; this draft does not independently verify that assessment.
\subsection{Sources}
\begin{itemize}
\item Local UnsolvedMath registry, permanent ID \(30005406\), title: Positive Characteristic Cohomology of Flag Variety Line Bundles. Collection: Oberwolfach Reports Open Problems. Exact source code: \texttt{catalogue entry}; the supplied original excerpt labels the question ``Problem 3''. The record gives the proposed year as 2023 and supplies no individual proposer.
\item \emph{Resolutions in Local Algebra and Singularity Theory} (2023), Oberwolfach Reports, pp.\ 327--395, DOI \texttt{10.4171/owr/2023/6}. Supplied report citation and literature-source label ``original OWR problem'': \url{https://doi.org/10.4171/owr/2023/6}. The extracted text includes a page artifact identifying p.\ 341. The record describes Oberwolfach Reports as documenting mathematical workshops and their open-problem discussions.
\item Supplied literature source, labeled ``Gao--Raicu, incidence correspondence'': \url{https://arxiv.org/abs/2210.03656}.
\item Supplied literature source, labeled ``cohomology on the incidence correspondence'': \url{https://arxiv.org/abs/2411.13450}.
\item Literature assessment date: \texttt{2026-08-22}. This is supplied dated curation, not independently verified. It labels the problem partially solved, states that the broad positive-characteristic flag-variety problem remains open, and reports complete vanishing and nonvanishing results for the incidence correspondence and recursive formulas for related flags. The supplied background explicitly requests independent verification of time-sensitive status and cited claims.
\item The registry labels the statement \texttt{corrected\_verified} and reports that the official report was checked directly and that the clean statement preserves the complete mathematical target while restoring definitions and removing extraction artifacts or source-dependent wording. These are supplied verification claims, not checks performed for this draft. The source-local bibliography entries [9, 8] for Kempf and [2] for Andersen are not included in the supplied record.
\end{itemize}
\end{document}
